\documentclass[12pt]{article}
\usepackage[utf8]{inputenc}
\usepackage[T1]{fontenc}
\usepackage[margin=1in]{geometry}
\usepackage{setspace}
\usepackage{xcolor}
\usepackage[colorlinks=true, urlcolor=blue!60!black, linkcolor=blue!60!black]{hyperref}
\setlength{\parindent}{0em}
\setlength{\parskip}{0.6em}
\renewcommand{\baselinestretch}{1.0}
\begin{document}
\begin{center}
{\LARGE \textit{Corfield v. Coryell}, 6 F. Cas. 546 (C.C.E.D. Pa. 1823)\par}
\vspace{0.5em}
{\large Case Brief --- Two-Page Edition\par}
\end{center}

\section{Memory Jogger}
Justice Bushrod Washington, on circuit, upheld New Jersey's ban on nonresident oyster harvesting: Article IV protects only ``fundamental'' rights, not a share in a state's common property (\href{https://press-pubs.uchicago.edu/founders/documents/a4_2_1s18.html}{Founders' Constitution}).

\section{Detailed Case Facts}
In 1821 New Jersey officers seized Corfield's vessel, the \textit{Hiram}, leased through Hand to John Keene, while Keene dredged for oysters in Maurice River Cove. A New Jersey statute barred nonresidents from gathering oysters in state waters aboard vessels not wholly owned by residents, on pain of forfeiture. Local magistrates condemned the vessel, and it was sold for \$10.

\section{Procedural History}
Corfield sued the prize master, Coryell, in trespass in federal circuit court. The jury returned a \$560 verdict subject to reserved questions of law; Justice Washington later ordered judgment for Coryell.

\section{Judicial Votes}
A single opinion by Circuit Justice Bushrod Washington; no concurrences or dissents.

\section{Holding}
The statute violates neither the Commerce, Privileges and Immunities, nor Admiralty Clause. Article IV protects rights fundamental in their nature, not a cotenancy in a state's oyster beds. Independently, a lessor out of possession cannot sue in trespass---so the constitutional rulings, and especially the catalogue of rights, are arguably dicta.

\section{Analysis of Opinions}
\textbf{Opinion of the court:} Under \textit{Gibbons v. Ogden}, Congress's power over navigation did not transfer a state's proprietary interest in its fisheries. Article IV does not entitle nonresidents to every local advantage; it secures fundamental rights such as travel and residence for trade or professional pursuits, habeas corpus, access to courts, property, equal taxation, and the franchise as regulated by state law. The state's citizens hold its fisheries as ``tenants in common,'' and opening them to all would risk exhausting a finite resource. The states never surrendered their fisheries through the admiralty grant, and Corfield, having leased his vessel, lacked possession to sue.

\section{Examples: Future Applications}
\textbf{Same-Side:} higher elk-hunting fees for nonresident recreational hunters are valid (\textit{Baldwin}, 1978). \textbf{Opposite-Side:} a prohibitive license fee on nonresident commercial shrimpers burdens a livelihood and is invalid (\textit{Toomer}, 1948). \textbf{Fence-Sitter:} a city's 40\% resident-hiring rule on city-funded projects implicates a fundamental calling but may be justified (\textit{Camden}, 1984).

\section{Critique}
Originalists divide over whether Washington described substantive natural rights later built into the Fourteenth Amendment (Barnett, Epstein) or a mere comity guarantee (Lash, Berger). Others fault his natural-law vagueness and the common-property fiction, which can mask protectionism.

\section{Key Quotations}
\begin{quote}
``Protection by the government; the enjoyment of life and liberty, with the right to acquire and possess property of every kind, and to pursue and obtain happiness and safety; subject nevertheless to such restraints as the government may justly prescribe for the general good of the whole.''
\end{quote}
\noindent --- Circuit Justice Washington, 6 F. Cas. at 551--52.

\end{document}
